\documentclass[11pt,a4paper]{article}

\usepackage[utf8]{inputenc}
\usepackage[T1]{fontenc}
\usepackage[french]{babel}
\usepackage[a4paper,margin=2.1cm]{geometry}
\usepackage{amsmath,amssymb,amsthm}
\usepackage{enumitem}
\usepackage{tabularx}
\usepackage{array}
\usepackage{xcolor}
\usepackage{lmodern}
\usepackage{fancyhdr}
\usepackage{lastpage}
\usepackage{hyperref}
\usepackage{multicol}
\usepackage{titlesec}
\usepackage{mathrsfs}

\hypersetup{
    colorlinks=true,
    linkcolor=black,
    urlcolor=blue
}

\pagestyle{fancy}
\fancyhf{}
\lhead{Terminale spécialité -- Mathématiques}
\rhead{DM -- Analyse de fonctions}
\cfoot{\thepage/\pageref{LastPage}}

\titleformat{\section}{\large\bfseries}{\thesection.}{0.5em}{}
\titleformat{\subsection}{\normalsize\bfseries}{\thesubsection.}{0.5em}{}

\setlist[itemize]{leftmargin=1.8em}
\setlist[enumerate]{leftmargin=2em}

\newcolumntype{Y}{>{\raggedright\arraybackslash}X}

\begin{document}

\begin{center}
    {\LARGE \textbf{Devoir Maison -- Mathématiques}}\\[0.3em]
    {\Large \textbf{Terminale spécialité}}\\[0.5em]
    {\large \textbf{Thème : Analyse de fonctions}}\\[0.8em]
\end{center}

\vspace{0.5em}

\noindent\textbf{Consignes générales :}
\begin{itemize}
    \item Les réponses doivent être \textbf{rédigées}, \textbf{justifiées} et organisées avec soin.
    \item Les raisonnements, les interprétations graphiques, les tableaux de variations et les conclusions seront particulièrement valorisés.
    \item Les résultats numériques ou algébriques non justifiés ne rapporteront qu'une partie des points.
    \item Le barème donne davantage de points aux questions longues, aux études globales de fonctions et aux raisonnements fins qu'aux applications directes du cours.
\end{itemize}

\vspace{0.8em}

\begin{center}
    \fbox{\parbox{0.93\linewidth}{
    \textbf{Répartition indicative du barème :} \\
    Problème 1 : \textbf{25 points} \hfill
    Problème 2 : \textbf{29 points} \hfill
    Problème 3 : \textbf{30 points} \\[0.2em]
    \hfill \textbf{Total : 84 points}
    }}
\end{center}

\section*{Problème 1 -- Étude complète d'une fonction logarithmique \hfill /25}

On considère la fonction \(f\) définie sur \(]0,+\infty[\) par
\[
f(x)=x-1-\ln(x).
\]

Ce problème est volontairement construit dans l'esprit des études classiques du bac : domaine, dérivée, variations, convexité, tangente, lecture graphique et inégalités.

\subsection*{1. Premières observations}
\begin{enumerate}
    \item Déterminer l'ensemble de définition de \(f\). \hfill \textbf{/1}
    \item Calculer \(f(1)\). \hfill \textbf{/1}
    \item Étudier la limite de \(f(x)\) lorsque \(x\to 0^+\). \hfill \textbf{/2}
    \item Étudier la limite de \(f(x)\) lorsque \(x\to +\infty\). \hfill \textbf{/1}
\end{enumerate}

\subsection*{2. Variations de la fonction}
\begin{enumerate}
    \item Calculer \(f'(x)\) pour tout \(x>0\). \hfill \textbf{/2}
    \item Étudier le signe de \(f'(x)\). \hfill \textbf{/2}
    \item En déduire le tableau de variations complet de \(f\). \hfill \textbf{/3}
    \item Montrer que \(f(x)\geq 0\) pour tout \(x>0\). \hfill \textbf{/2}
\end{enumerate}

\subsection*{3. Étude de la convexité}
\begin{enumerate}
    \item Calculer \(f''(x)\). \hfill \textbf{/1}
    \item Étudier le signe de \(f''(x)\). \hfill \textbf{/1}
    \item En déduire la convexité de la fonction \(f\). \hfill \textbf{/1}
    \item Donner une interprétation graphique de ce résultat. \hfill \textbf{/1}
\end{enumerate}

\subsection*{4. Tangente et comparaison locale}
\begin{enumerate}
    \item Déterminer une équation de la tangente à la courbe de \(f\) au point d'abscisse \(1\). \hfill \textbf{/2}
    \item Étudier la position relative de la courbe de \(f\) et de cette tangente. \hfill \textbf{/3}
    \item Retrouver ainsi une inégalité simple reliant \(\ln(x)\) et \(x-1\). \hfill \textbf{/2}
\end{enumerate}

\subsection*{5. Regard critique}
\begin{enumerate}
    \item Un élève affirme : \og comme \(f(1)=0\), la fonction change forcément de signe en \(1\). \fg \\
    Dire si cette affirmation est vraie ou fausse, puis justifier soigneusement. \hfill \textbf{/2}
    \item Expliquer pourquoi ce problème permet de comprendre géométriquement que la fonction logarithme est en dessous de certaines de ses tangentes. \hfill \textbf{/1}
\end{enumerate}

\vspace{1em}

\section*{Problème 2 -- Une fonction rationnelle-exponentielle pour modéliser un rendement \hfill /29}

On considère la fonction \(g\) définie sur \([0,+\infty[\) par
\[
g(x)=\frac{x}{x+1}\,e^{-x}.
\]

Cette fonction modélise de façon simplifiée un rendement : le facteur \(\dfrac{x}{x+1}\) augmente avec \(x\), tandis que le facteur \(e^{-x}\) traduit une perte progressive.

\subsection*{1. Étude préliminaire}
\begin{enumerate}
    \item Déterminer l'ensemble de définition de \(g\). \hfill \textbf{/1}
    \item Calculer \(g(0)\). \hfill \textbf{/1}
    \item Déterminer la limite de \(g(x)\) lorsque \(x\to +\infty\). \hfill \textbf{/2}
    \item Interpréter graphiquement les résultats précédents. \hfill \textbf{/1}
\end{enumerate}

\subsection*{2. Dérivation et variations}
\begin{enumerate}
    \item Montrer que, pour tout \(x\geq 0\),
    \[
    g'(x)=\frac{e^{-x}}{(x+1)^2}\bigl(1-x-x^2\bigr).
    \]
    \hfill \textbf{/4}
    \item Étudier le signe de \(g'(x)\) sur \([0,+\infty[\). \hfill \textbf{/3}
    \item En déduire le tableau de variations complet de \(g\). \hfill \textbf{/3}
    \item Montrer que l'équation \(1-x-x^2=0\) admet une unique solution positive \(\alpha\), que l'on exprimera exactement. \hfill \textbf{/2}
    \item En déduire que \(g\) admet un maximum sur \([0,+\infty[\), atteint pour \(x=\alpha\). \hfill \textbf{/1}
\end{enumerate}

\subsection*{3. Exploitation du maximum}
\begin{enumerate}
    \item Calculer la valeur exacte de \(\alpha\). \hfill \textbf{/1}
    \item Donner une valeur approchée au centième de \(\alpha\). \hfill \textbf{/1}
    \item Donner une valeur approchée au centième de \(g(\alpha)\). \hfill \textbf{/2}
    \item Interpréter concrètement le rôle de \(\alpha\) dans le modèle. \hfill \textbf{/2}
\end{enumerate}

\subsection*{4. Une équation associée}
On considère l'équation
\[
g(x)=0,1
\]
sur l'intervalle \([0,+\infty[\).

\begin{enumerate}
    \item Expliquer, sans résoudre explicitement l'équation, pourquoi elle admet soit zéro, soit une, soit deux solutions sur \([0,+\infty[\). \hfill \textbf{/2}
    \item En utilisant les variations de \(g\), expliquer pourquoi cette équation admet en fait deux solutions. \hfill \textbf{/3}
    \item Donner un encadrement de chacune des solutions d'amplitude \(1\), à l'aide d'un simple tableau de valeurs si nécessaire. \hfill \textbf{/2}
\end{enumerate}

\subsection*{5. Prise de recul}
\begin{enumerate}
    \item Expliquer pourquoi la présence du facteur exponentiel modifie profondément le comportement global de la fonction par rapport à la seule fonction \(\dfrac{x}{x+1}\). \hfill \textbf{/2}
    \item Sans faire de calcul complet, dire ce que deviendrait qualitativement l'étude si on remplaçait \(e^{-x}\) par \(e^{-2x}\). \hfill \textbf{/2}
\end{enumerate}

\vspace{1em}

\textbf{Problème 3 -- Une équation fonctionnelle autour des symétries d'une courbe \hfill /30}

On considère une fonction \(f\) définie sur \(\mathbb{R}\) par :
\[
f(x)=\frac{1}{1+e^{-x}}.
\]

On cherche à étudier cette fonction puis à mettre en évidence une relation remarquable reliant \(f(x)\) et \(f(-x)\).

\subsection*{1. Premières propriétés}
\begin{enumerate}
    \item Montrer que \(f\) est définie sur \(\mathbb{R}\). \hfill \textbf{/1}
    \item Calculer \(f(0)\). \hfill \textbf{/1}
    \item Déterminer les limites de \(f(x)\) lorsque \(x\to -\infty\) puis lorsque \(x\to +\infty\). \hfill \textbf{/3}
    \item Interpréter graphiquement ces résultats. \hfill \textbf{/1}
\end{enumerate}

\subsection*{2. Variations de la fonction}
\begin{enumerate}
    \item Montrer que, pour tout réel \(x\),
    \[
    f'(x)=\frac{e^{-x}}{(1+e^{-x})^2}.
    \]
    \hfill \textbf{/3}
    \item Étudier le signe de \(f'(x)\). \hfill \textbf{/2}
    \item En déduire le tableau de variations de \(f\) sur \(\mathbb{R}\). \hfill \textbf{/2}
\end{enumerate}

\subsection*{3. Une équation fonctionnelle}
\begin{enumerate}
    \item Calculer \(f(-x)\). \hfill \textbf{/2}
    \item Montrer que, pour tout réel \(x\),
    \[
    f(x)+f(-x)=1.
    \]
    \hfill \textbf{/4}
    \item Expliquer pourquoi cette relation est appelée une \textbf{équation fonctionnelle}. \hfill \textbf{/1}
\end{enumerate}

\subsection*{4. Interprétation graphique}
\begin{enumerate}
    \item Soit \(M(x,f(x))\) un point de la courbe représentative de \(f\). Montrer que le point
    \[
    M'(-x,1-f(x))
    \]
    appartient aussi à cette courbe. \hfill \textbf{/3}
    \item En déduire que la courbe de \(f\) admet un centre de symétrie dont on précisera les coordonnées. \hfill \textbf{/3}
    \item Vérifier ce résultat sur le point correspondant à \(x=0\). \hfill \textbf{/1}
\end{enumerate}

\subsection*{5. Une autre écriture utile}
\begin{enumerate}
    \item Montrer que, pour tout réel \(x\),
    \[
    1-f(x)=\frac{e^{-x}}{1+e^{-x}}.
    \]
    \hfill \textbf{/2}
    \item Montrer alors que
    \[
    f'(x)=f(x)\bigl(1-f(x)\bigr).
    \]
    \hfill \textbf{/3}
    \item Expliquer en quoi cette écriture permet de mieux comprendre le signe de \(f'(x)\). \hfill \textbf{/1}
\end{enumerate}

\subsection*{6. Question de recul}
\begin{enumerate}
    \item Un élève affirme : \og comme \(f(x)+f(-x)=1\), la fonction \(f\) est paire \fg. Dire si cette affirmation est vraie ou fausse et justifier soigneusement. \hfill \textbf{/2}
    \item Expliquer la différence entre :
    \begin{itemize}
        \item une symétrie par rapport à l'axe des ordonnées ;
        \item une symétrie centrale.
    \end{itemize}
    \hfill \textbf{/1}
\end{enumerate}

\vspace{1em}

\section*{Barème détaillé récapitulatif}

\renewcommand{\arraystretch}{1.25}
\begin{tabularx}{\linewidth}{|Y|c|}
\hline
\textbf{Éléments évalués} & \textbf{Points} \\
\hline
\textbf{Problème 1} : étude classique de fonction, variations, convexité, tangente, interprétation d'inégalités & \textbf{25} \\
\hline
\textbf{Problème 2} : dérivation plus technique, maximum, exploitation d'un modèle, existence de solutions, lecture qualitative & \textbf{29} \\
\hline
\textbf{Problème 3} : fonction définie par intégrale, croissance, convexité, majoration, limite, lien avec une primitive connue & \textbf{30} \\
\hline
\textbf{Total} & \textbf{84} \\
\hline
\end{tabularx}

\vspace{1em}

\noindent\textbf{Remarque sur la philosophie du barème :}
\begin{itemize}
    \item Les questions de \textbf{cours direct} ou de calcul immédiat rapportent peu.
    \item Les questions demandant une \textbf{étude complète}, un \textbf{raisonnement global}, une \textbf{interprétation graphique} ou une \textbf{prise de recul} rapportent davantage.
    \item Une réponse juste mais non justifiée ne rapporte qu'une partie des points.
\end{itemize}

\vspace{1em}

\begin{center}
    \fbox{\parbox{0.92\linewidth}{
    \textbf{Conseil à l'élève :} Pour réussir ce DM, il faut aller au-delà du calcul. On attend une vraie compréhension des variations, de la convexité, du rôle des dérivées et du sens des objets étudiés.
    }}
\end{center}

\end{document}